\subsection{Sobolev 空间的实插值}
\label{sec:ch14-real-interpolation-sobolev-spaces}
\setcounter{thmcount}{0}
\setcounter{equation}{0}

最基本的插值结果 Riesz 凸性定理把 Lebesgue 空间联系起来. 在我们的设定中, 它就是
\MBSourceItem{1}
\begin{equation}
\MathBookFormulaID{formula-ch14-riesz-convexity-lebesgue-interpolation}
\label{eq:ch14-riesz-convexity-lebesgue-interpolation}
L^p(\Omega)=[L^1(\Omega),L^\infty(\Omega)]_{1-1/p,p}
\qquad\forall,1<p<\infty
\end{equation}
(参见 Bergh 与 Löfstrom 1976 以及 Stein 与 Weiss 1971). 这也可推广到 Sobolev 空间:
\MathBookSourcePage{382}
\MBSourceItem{2}
\begin{equation}
\MathBookFormulaID{formula-ch14-sobolev-p-interpolation}
\label{eq:ch14-sobolev-p-interpolation}
W_p^k(\Omega)=[W_1^k(\Omega),W_\infty^k(\Omega)]_{1-1/p,p}
\qquad\forall,1<p<\infty,
\end{equation}
其中范数等价, 参见 DeVore 与 Scherer (1979). 到目前为止, 我们只在固定 $k$ 时对变量 $p$ 作了插值. 现在固定 $p$ 而改变 $k$.

分数阶 Sobolev 空间的另一种定义由插值得到. 下面证明这一定义与此前定义等价.
\MBSourceItem{3}
\begin{Theorem}
\label{thm:ch14-fractional-sobolev-real-interpolation}
设 $0<s<1$. 若 $\Omega$ 具有 Lipschitz 边界, 则
\[
\MathBookFormulaID{formula-ch14-fractional-sobolev-interpolation-space}
W_p^{k+s}(\Omega)=[W_p^k(\Omega),W_p^{k+1}(\Omega)]_{s,p},
\]
且范数等价.
\end{Theorem}

\begin{Proof}
当 $\Omega=\mathbb R^n$ 时, 已知两种定义给出相同的空间和等价的范数. 例如, Adams (1975) 使用一种与前述 $K$ 方法等价的插值方法(迹方法)证明了这一点(Bergh 与 Löfstrom 1976). 另见 Bergh 与 Löfstrom (1976) 的定理 6.4.5 第 (4) 项以及第 6 章习题 7.

由前述 Stein 延拓结果定理~\ref{thm:ch01-extension-mapping}, 对所有 $k$ 都存在算子 $E_S:W_p^k(\Omega)\to W_p^k(\mathbb R^n)$. 对该算子作插值, 得
\[
\MathBookFormulaID{formula-ch14-extension-interpolation-forward-bound}
\begin{aligned}
\|u\|_{W_p^{k+s}(\Omega)}
&=\|E_Su\|_{W_p^{k+s}(\Omega)}\\
&\leq\|E_Su\|_{W_p^{k+s}(\mathbb R^n)} &&\text{由~\eqref{eq:ch14-intrinsic-fractional-sobolev-norm}}\\
&\leq C\|E_Su\|_{[W_p^k(\mathbb R^n),W_p^{k+1}(\mathbb R^n)]_{s,p}} &&\text{使用 $\mathbb R^n$ 情形}\\
&\leq C\|u\|_{[W_p^k(\Omega),W_p^{k+1}(\Omega)]_{s,p}} &&\text{由命题~\ref{prop:ch14-operator-interpolation}}.
\end{aligned}
\]
注意, 此论证中必须知道算子 $E_S$ 同时定义在 $W_p^k(\Omega)$ 和 $W_p^{k+1}(\Omega)$ 上, 否则不能应用算子插值.

反过来, Grisvard (1985) 指出, 若 $\Omega$ 是 Lipschitz 区域, 则存在延拓算子 $E_G$, 它把 $W_p^{k+s}(\Omega)$ 映到 $W_p^{k+s}(\mathbb R^n)$, 且 $(E_Gu)|_\Omega=u$; 这里分数阶 Sobolev 范数按~\eqref{eq:ch14-intrinsic-fractional-sobolev-norm} 定义. 证明见 DeVore 与 Sharpley (1993). 因而
\[
\MathBookFormulaID{formula-ch14-extension-interpolation-reverse-bound}
\begin{aligned}
\|u\|_{[W_p^k(\Omega),W_p^{k+1}(\Omega)]_{s,p}}
&=\|E_Gu\|_{[W_p^k(\Omega),W_p^{k+1}(\Omega)]_{s,p}}\\
&\leq\|E_Gu\|_{[W_p^k(\mathbb R^n),W_p^{k+1}(\mathbb R^n)]_{s,p}}
&&\text{由~\eqref{eq:ch14-domain-restriction-interpolation-inclusion} 与习题~\ref{exr:ch14-02}}\\
&\leq C\|E_Gu\|_{W_p^{k+s}(\mathbb R^n)} &&\text{使用 $\mathbb R^n$ 情形}\\
&\leq C\|u\|_{W_p^{k+s}(\Omega)}.
\end{aligned}
\]
这里允许 $E_G$ 依赖于 $s$. 因而在 Lipschitz 区域的情形, 内蕴范数~\eqref{eq:ch14-intrinsic-fractional-sobolev-norm} 与实插值得到的范数等价.
\end{Proof}

\MathBookSourcePage{383}
对更复杂的区域, 定理~\ref{thm:ch14-fractional-sobolev-real-interpolation} 中考虑的空间未必等价, 这一点并不明确.

除非 $p=2$, 空间 $W_p^{k+1}(\Omega)$ 与 $[W_p^k(\Omega),W_p^{k+2}(\Omega)]_{1/2,p}$ 并不相同, 不过已知它们非常接近(Bergh 与 Löfstrom 1976), 即
\[
\MathBookFormulaID{formula-ch14-nearby-second-order-interpolation-spaces}
[W_p^k(\Omega),W_p^{k+2}(\Omega)]_{1/2,1}
\subset W_p^{k+1}(\Omega)
\subset [W_p^k(\Omega),W_p^{k+2}(\Omega)]_{1/2,\infty}.
\]
然而, 定理~\ref{thm:ch14-fractional-sobolev-real-interpolation} 的证明还表明
\MBSourceItem{4}
\begin{equation}
\MathBookFormulaID{formula-ch14-general-sobolev-real-interpolation}
\label{eq:ch14-general-sobolev-real-interpolation}
[W_p^m(\Omega),W_p^k(\Omega)]_{\theta,p}
=W_p^{(1-\theta)m+\theta k}(\Omega),
\end{equation}
只要 $(1-\theta)m+\theta k$ 不是整数. 此外, 重申定理~\eqref{eq:ch14-reiteration-theorem} 说明, 即使 $m$ 和 $k$ 不是整数, ~\eqref{eq:ch14-general-sobolev-real-interpolation} 仍成立.

与整数情形一样, 可以用对偶性定义负分数指标 Sobolev 空间, 即 $W_p^{-s}(\Omega):=W_{p'}^s(\Omega)'$. 对偶定理~\eqref{eq:ch14-interpolation-duality} 说明, ~\eqref{eq:ch14-general-sobolev-real-interpolation} 对负指标也成立.

强调当 $p=2$ 时, 关系~\eqref{eq:ch14-general-sobolev-real-interpolation} 无限制成立, 即
\MBSourceItem{5}
\begin{equation}
\MathBookFormulaID{formula-ch14-hilbert-sobolev-interpolation}
\label{eq:ch14-hilbert-sobolev-interpolation}
[H^m(\Omega),H^k(\Omega)]_{\theta,2}=H^{(1-\theta)m+\theta k}(\Omega).
\end{equation}
一旦知道 $\Omega=\mathbb R^n$ 时结论成立, 定理~\ref{thm:ch14-fractional-sobolev-real-interpolation} 的证明便给出该结果. 后者可利用 $H^s(\mathbb R^n)$ 范数的另一种等价定义证明, 即
\MBSourceItem{6}
\begin{equation}
\MathBookFormulaID{formula-ch14-fourier-sobolev-norm}
\label{eq:ch14-fourier-sobolev-norm}
\|u\|_{H^s(\mathbb R^n),F}
:=\left(\int_{\mathbb R^n}(1+|\xi|^2)^s|\widehat u(\xi)|^2\,d\xi\right)^{1/2},
\end{equation}
其中 $\widehat u$ 表示 $u$ 的 Fourier 变换. 以此为中介(参见习题~\ref{exr:ch14-06} 和~\ref{exr:ch14-07}), 即可完成等价性的证明.

最后注意, ~\eqref{eq:ch14-hilbert-sobolev-interpolation} 实际上不要求 $m$ 与 $k$ 同号. 首先有
\[
\MathBookFormulaID{formula-ch14-symmetric-index-interpolation}
[H^{-k}(\Omega),H^k(\Omega)]_{1/2,2}=L^2(\Omega)=H^0(\Omega)
\]
(参见 Bergh 与 Löfstrom 1976). 使用重申定理~\eqref{eq:ch14-reiteration-theorem}, 得到下述结果.
\MBSourceItem{7}
\begin{Theorem}
\label{thm:ch14-arbitrary-real-sobolev-interpolation}
若 $\Omega$ 具有 Lipschitz 边界, 则~\eqref{eq:ch14-hilbert-sobolev-interpolation} 对任意实数 $m$ 和 $k$ 成立, 且范数等价.
\end{Theorem}
